% Gerado por avaliacao/relatorio.py — nao editar a mao; rode novamente a avaliacao.
\begin{table}[htbp]
\caption{Desempenho da recuperação por dispositivo}\label{tab:recuperacao}
\centering\footnotesize
\setlength{\tabcolsep}{4pt}
\begin{tabular}{lrrrrrrr}
\toprule
Recuperador & P@1 & R@3 & R@6 (IC 95\%) & R@10 & Acerto@6 & MRR & R@6 artigo \\
\midrule
hash & 12,5\% & 13,5\% & 15,4\% (7--26) & 15,4\% & 28,1\% & 0,188 & 43,0\% \\
spaCy & 9,4\% & 9,4\% & 16,9\% (7--29) & 19,0\% & 28,1\% & 0,154 & 33,9\% \\
BGE-M3 & 21,9\% & 21,3\% & 24,2\% (13--36) & 34,6\% & 43,8\% & 0,312 & 46,6\% \\
BM25 & 25,0\% & 31,5\% & 36,5\% (23--50) & 40,4\% & 56,2\% & 0,374 & 71,1\% \\
Híbrido (BGE-M3 + BM25) & 31,2\% & 28,9\% & 39,6\% (27--53) & 47,9\% & 68,8\% & 0,427 & 67,5\% \\
\bottomrule
\end{tabular}
\par\smallskip
\begin{minipage}{\linewidth}{\small \textbf{Fonte:} Elaborado pelo autor (2026), a partir de \texttt{avaliacao/resultados/resumo.json}.}\par
{\small \textbf{Nota:} 32 demandas, 86 dispositivos relevantes no gabarito, 1.734 dispositivos na base. IC: intervalo de confiança por bootstrap (2.000 reamostragens). R@6 artigo: acerto quando o dispositivo recuperado pertence ao mesmo artigo do esperado.}\end{minipage}
\end{table}
